% Related Work

\section{Related Work}
\label{sec:related}

\paragraph{Step-wise and verifier GUI agents.}
VLM-based mobile GUI agents usually follow an observe--reason--act loop: each
screen update is converted into a new model decision, and the system issues one
concrete GUI action before observing again.  \textsc{V-Droid}~\cite{dai2025advancing}
improves this loop by shifting the LLM role from action generation to action
verification over a discrete action set, reporting 4.3\,s per step and 59.5\%
success on AndroidWorld.  \textsc{VeriSafe}~\cite{verisafe} similarly focuses
on action safety by checking LLM-generated actions against logic constraints
derived from the user's goal.  These systems improve the cost or correctness of
individual steps, but repeated routines still pass through online decision or
verification for each interaction.  OmniFlow is complementary: a step-wise or
verifier-based agent remains the fallback, while matched repeated Functions are
served as executable Function cache hits.

\paragraph{Memory-based and programmatic GUI agents.}
\textsc{AutoDroid}~\cite{wen2024autodroid},
\textsc{MobileGPT}~\cite{lee2024mobilegpt}, and
\textsc{AppAgent}~\cite{zhang2025appagent} reduce repeated GUI-agent work by
building app- or task-specific memory from exploration and prior executions.
AndroidWorld's M3A baseline~\cite{rawles2025androidworld} also exposes the
benefit of carrying previous actions and reasoning into future attempts.
These systems reduce per-goal LLM cost when a warm-start match is found, but
their memory is primarily used to guide online reasoning or replay app-specific
subtasks.  Reuse is limited by per-app cold-start exploration and
UI-attribute-based action grounding: XML selectors break on WebViews,
custom-rendered feeds, and cross-device layout variation, which forces repair,
new exploration, or online LLM calls.
\textsc{ActionEngine}~\cite{zhong2026actionengine} and
\textsc{AutoDroid-V2}~\cite{autodroid-v2}, as well as
\textsc{AgentProg}~\cite{tian2025agentprog}, move toward programmatic
execution, building app-level state knowledge or generated scripts to
reduce online planning.
Their reuse unit is a broad app model or complete task program that
requires substantial upfront exploration and assumes a stable live
app structure.
\sysname takes a narrower boundary: it treats a verified source run as the
asset, caches a Function-level action list, and focuses on transferring each
recorded source action to the target GUI instance.


\paragraph{GUI grounding.}
Visual grounding models for GUI, such as
\textsc{SeeClick}~\cite{cheng2024seeclick},
\textsc{UGround}~\cite{gou2024uground}, and
\textsc{ShowUI}~\cite{lin2025showui},
train VLMs to map a natural-language element description to a screen
coordinate or bounding box.
These models optimize for the \emph{zero-prior} setting: given only a
task instruction and the current screenshot, they locate the target
element without any recorded demonstration.
This makes them general-purpose and immediately applicable to novel apps
and screens.
OmniFlow addresses the complementary setting: a human demonstration
already exists, and the goal is to rebind its recorded actions to the
current screen without a model call.
GUI grounding models could in principle serve as the transfer oracle in
OmniFlow's execution loop, but their per-step inference cost reintroduces
the model overhead that OmniFlow is designed to eliminate.
Feature-engineered element vectors let \textsc{Flow} transfer recorded
targets through structural similarity alone, reserving model calls for
genuinely ambiguous or novel states.
 